\section{Introduction}\label{sec:introduction}

For a closed connected smooth manifold $Q$, write $\pi:T^*Q\to Q$ for
the cotangent projection and
$\lambda_{(q,p)}(v)=p(d\pi(v))$ for the canonical one-form. A smooth
submanifold $L\subset T^*Q$ is Lagrangian if it has dimension $\dim Q$
and $d\lambda$ vanishes on it; it is exact if $\lambda|_L=df$ for a
globally defined function $f:L\to\R$. Here closed means compact without
boundary. Arnold's nearby Lagrangian conjecture asks whether every closed
exact embedded Lagrangian in $T^*Q$ is Hamiltonian isotopic to the zero
section. This is the modern formulation used, for example, in
\cite{Abouzaid2012,Kragh2013}; Arnold's discussion of Lagrangian knots
provides the broader historical setting \cite{Arnold1986}.

Throughout this paper our conventions are
\[
  \omega=d\lambda,\qquad \iota_{X_{H_t}}\omega=-dH_t.
\]
A Hamiltonian isotopy has a smooth time-dependent Hamiltonian supported
in one fixed compact subset of $T^*Q$. For compact Lagrangians this
entails no restriction: multiplying a Hamiltonian by a cutoff equal to
one on a neighborhood of the compact isotopy trace preserves its flow
along that trace.

\begin{theorem}\label{thm:main}
For some sufficiently large even integer $N$, the manifold
$Q=S^9\times S^{N-1}$ admits a closed exact smoothly embedded Lagrangian
$L\subset T^*Q$ such that $L$ is diffeomorphic to $Q$ and no compactly
supported Hamiltonian isotopy carries the zero section to $L$.
\end{theorem}

The theorem resolves the unrestricted nearby Lagrangian conjecture
negatively. Both the base and the Lagrangian in this example are
connected and simply connected. The obstruction concerns how $L$ lies
in the cotangent bundle, even though its abstract smooth manifold agrees
with the base.

\subsection{Context and predecessors}

Classical intersection theorems of Hofer and Laudenbach--Sikorav constrain
Hamiltonian deformations of the zero section
\cite{Hofer1985,LaudenbachSikorav1985}. Generating functions provide one
route to such results: their critical points encode intersections, so
finite-dimensional critical-point theory can detect symplectic phenomena.
The construction and transport of these functions under Hamiltonian
isotopy were developed by Chaperon, Laudenbach and Sikorav
\cite{Chaperon1984,LaudenbachSikorav1985,Sikorav1987}; Eliashberg and Gromov give a systematic
account, including Chekanov's contact-isotopy lifting theorem
\cite[Theorem~4.1.1]{EG1998}. Positive isotopy results in low
dimension include Hind's theorem for Lagrangian spheres in $T^*S^2$
\cite[Theorem~1]{Hind2012} and the cotangent-torus theorem of Dimitroglou
Rizell, Goodman and Ivrii \cite[Theorem~B]{DGI2016}. The former, together
with the projection theorem recalled next and the classification of
closed surfaces, gives the closed exact case over $S^2$.

The topology of the cotangent projection has become increasingly rigid.
Fukaya, Seidel and Smith, and independently Nadler, established
cohomological and categorical constraints under additional assumptions
on the base, Maslov class and spin structures \cite{FSS2008,Nadler2009}.
Abouzaid proved that the projection is a homotopy equivalence when the
Maslov class vanishes \cite{Abouzaid2012}. Kragh's
parametrized spectra, together with Abouzaid's appendix, removed that
hypothesis \cite[Corollary~1.2 and Theorem~E.2]{Kragh2013}.
Abouzaid and Kragh subsequently proved that the projection is a simple
homotopy equivalence \cite[Theorem~1]{AK2018}. These results concern the
topology of the projection and do not identify the ambient Hamiltonian
isotopy class.

More recent work reaches finer topological information. Abouzaid,
Courte, Guillermou and Kragh constructed twisted generating functions
and proved that the stable Lagrangian Gauss map, which records the
stabilized tangent Lagrangian planes, induces zero on homotopy groups
\cite{ACGK2025}. The normal invariant of a homotopy equivalence to $Q$
lies in $[Q,G/O]$, where $G/O$ is the homotopy fiber of the map from
stable real vector bundles to stable spherical fibrations. It measures
part of the obstruction to making the equivalence a diffeomorphism.
Abouzaid, \mbox{\'{A}lvarez-Gavela},
Courte and Kragh used the twisted derivative on tube spaces to constrain
the normal invariant of the projection, in particular proving that it
is $2$-torsion \cite[Corollary~1.2]{AACK2025}. Porcelli and Smith's
Floer-homotopy approach places this normal invariant in the image of the
stable Hopf action $B(G/O)\to G/O$ \cite[Theorem~6]{PorcelliSmith2026}.
These are further necessary conditions on a nearby
Lagrangian, distinct from its ambient isotopy class.

Smooth tube spaces connect generating functions with Waldhausen's
algebraic $K$-theory of spaces \cite{Waldhausen1982,WJR2013}.
\mbox{\'{A}lvarez-Gavela}, Igusa and Sullivan developed tube torsion for
Legendrians and constructed Legendrians whose tube torsion is nontrivial
although their base projections are homotopic to diffeomorphisms
\cite[Theorem~1.15]{AIS2026}. They also identify stable tube homotopy
classes as a finer invariant underlying the real cohomological torsion
\cite[Remark~3.32]{AIS2026}. Their examples lie in the jet space
$J^1Q=T^*Q\times\R$; forgetting the last coordinate can introduce double
points. The present construction uses a particular nonzero stable tube
class and proves that every possible pair of coincident cotangent
branches is separated.

Courte and Porcelli constrain the parametrized Whitehead torsion of
families of nearby Lagrangians, proving weak vanishing when the base has
Euler characteristic zero \cite[Theorems~1.2 and~1.5, Corollary~1.6]{CP2025}.
Their parameters index whole Lagrangians in a fixed cotangent bundle.
Below, $S^9$ instead indexes the functions at infinity of one generating
family. These two families, and the invariants attached to them, have
different roles. The one-critical-function and cobordism arguments are
also related to Kragh's study of generating functions and the
Hatcher--Waldhausen map \cite{Kragh2026}. We prove the required statement
for our varying homogeneous data at infinity.

\subsection{The construction and its obstruction}

We describe the proof in the language of generating families. For a
smooth function $F(q,w)$ on $Q\times\R^r$, the equations
\[
 d_wF(q,w)=0,\qquad (q,w)\longmapsto(q,d_qF(q,w))
\]
define an exact Lagrangian immersion when the full vertical-gradient
equation is regular. This means regularity with respect to both $q$ and
$w$, not invertibility of every fiber Hessian. An embedding requires the
additional global injectivity argument supplied in
Section~\ref{sec:embedding}.

Our generating families agree outside a fixed vertical ball with
degree-two homogeneous functions $g_b$, where $b\in S^9$. The negative
domain of $g_b$ on the unit sphere is a smooth tube: a domain isotopic
to the negative domain of a nondegenerate quadratic form. Adding positive and
negative quadratic variables defines the stable smooth tube space $\T$.
The homotopy type of each negative domain is a sphere, so such a family
also determines a stable spherical fibration, classified by a map to
$BG$. Section~\ref{sec:tubes} defines these models and their map
$\gamma:\T\to BG$. Waldhausen's tube fibration, together with Rognes's
low-degree smooth Whitehead group calculation
\cite{Waldhausen1982,Rognes2002,Rognes2014}, yields
\[
 0\ne a\in\pi_9\T,\qquad \gamma_*a=0.
\]
Thus the smooth tube family remains nontrivial after its associated
stable spherical fibration has become trivial.

Sections~\ref{sec:morse} and~\ref{sec:obstruction} prove that these
infinity data cannot occur in a family with exactly one nondegenerate
critical point in each fiber. Morse theory first gives a fiber homotopy
equivalence between the sphere of negative eigendirections at that point
and the negative domain at infinity. After the local quadratic model has been made constant, the
regular zero level between an inner and an outer sphere forms a family
of $h$-cobordisms. Thickening this family by one interval and extending
it into a fixed cylinder produces a cobordism with a product complement.
Stability and the classification of smooth families then give a product
for the whole original family. This permits the negative domains to be
deformed to a fixed tube, contradicting the choice of $a$.

Section~\ref{sec:embedding} realizes $a$ using the extra sphere factor
in $Q=S^9\times S^{N-1}$. A positive scale $R(b)$ and a linear coupling
to $v\in S^{N-1}$ arrange that the entire critical locus satisfies
\[
 \nabla g_b(w)=R(b)v.
\]
Homogeneity parametrizes this locus by $S^9\times S^{N-1}$. The sphere
and base momenta are the two components $d_vF$ and $d_bF$ of the generated
covector. Equality of the sphere momenta confines the difference between two putative
coincident branches to the line spanned by $v$. Varying $R$ over $S^9$
makes the base momenta separate the remaining pairs. Near the critical
points of that variation, the tube family has been chosen quadratic,
and the critical equation already has a unique solution.

Finally, Section~\ref{sec:transport} follows the generating family
through a hypothetical Hamiltonian isotopy to the zero section. Its
direct small-step construction adds fixed split quadratic forms and
preserves the stable class of the homogeneous data. At the zero section,
the critical-locus projection would be a diffeomorphism, giving a unique
nondegenerate critical point over each base point. Restricting to a fixed
$S^9$ slice then gives precisely the family excluded by
Proposition~\ref{prop:one-critical}, completing the contradiction.
